\section{跨海接收之后}

郑经1681年去世，东宁的继承冲突使郑克塽最终掌权（EQ-1681-ZHENG-JING-DEATH）。死期与继承结果可由清廷情报、郑氏后出材料和荷兰通信互校；宫廷政变的每一步不能因后出的叙述而获得同等确定性。三藩战争结束与东宁的不稳，给北京重新估算跨海行动提供了窗口。康熙二十一年任命施琅统率福建水师（EQ-1682-SHI-LANG），任命与筹备见于实录和军务奏折；施琅的战略意图及朝中争论须区别奏折、朱批与后出的传记，不能把成功后的理由当成当时唯一的选择。

康熙二十二年六月，施琅水师在澎湖击败郑氏舰队（EQ-1683-PENGHU）。胜败与其战略后果可以确立；舰船、兵员、死伤数字多出自胜方战报，须同郑氏叙事及海域材料交叉约束。季风、福建补给和航道知识，解释了行动为何能在那个时点集中展开，也提醒我们水师不是从地图上直线抵达台湾。八月郑克塽投降（EQ-1683-TAIWAN-SURRENDER），清廷取得政治接收权，但投降条件、官兵安置和移民、原住民社群的变化属于不同层次的过程，不能压缩为“台湾归入清朝”一项结果。

1684年设置台湾府及诸县、隶属福建，同时调整迁界与海禁，准许居民复界并逐步恢复海贸（EQ-1684-TAIWAN-PREFECTURE、EQ-1684-SEA-BAN-LIFT）。实录、地理志和福建档案能够证明制度设置和政策转向；府县与驻防的建立不等于全岛有效治理，复界诏令也不等于每一港口在同一时间开放。平原、沿海、内山与原住民社群同官府的接触程度并不相同。1721年朱一贵起事、一度进入府城，清廷继而派兵恢复控制（EQ-1721-ZHU-YIGUI）。军报可确立起事、府城易手和镇压的序列；官方“逆匪”分类不能代替参与者的社会构成或赋役诉求。地方志、督抚奏折和契约文书提示垦殖、租佃、治安及征收摩擦，却也只提供局部证词。跨海设治的脆弱性，正由这场起事暴露出来。


\subsection{事件链与材料限度}
东宁、台湾与清廷之间的转换须放在继承、航运与地方安置的不同节奏中阅读。

\paragraph{康熙二十年（1681年）：郑经去世，东宁发生继承冲突，郑克塽最终掌权}
\eventanchor{EQ-1681-ZHENG-JING-DEATH}{郑经去世，东宁发生继承冲突，郑克塽最终掌权}。本节点叙述该事件本身。东宁郑氏政权在京师与多地军政线路相接的尺度的处境首先受郑经长期在大陆作战并承受台湾财政压力，其死亡使宗族、将领和文官围绕继承重组制约。《清圣祖实录》康熙二十年台湾情报；《清史稿·郑成功传附郑经》；《台湾外纪》与荷兰商馆残存通信留下可核的线索；可辨的后果是东宁短期政治不稳，清廷得以评估以水师进攻台湾的风险和机会。原有置信注记：郑经死期和继承结果可靠；宫廷政变过程主要依郑氏后出材料与清方情报

\paragraph{康熙二十年（1681年）：郑经去世，东宁发生继承冲突，郑克塽最终掌权}
\eventanchor{EQ-1681-ZHENG-JING-DEATH-AFTERMATH}{郑经去世，东宁发生继承冲突，郑克塽最终掌权} 置于京师与多地军政线路相接的尺度中阅读，本节点追踪「郑经去世，东宁发生继承冲突，郑克塽最终掌权」之后可见的余波。郑经长期在大陆作战并承受台湾财政压力，其死亡使宗族、将领和文官围绕继承重组构成行动者东宁郑氏政权必须面对的条件。取证从《清圣祖实录》康熙二十年台湾情报；《清史稿·郑成功传附郑经》；《台湾外纪》与荷兰商馆残存通信出发，因而只能把变化谨慎地连到东宁短期政治不稳，清廷得以评估以水师进攻台湾的风险和机会。原有置信注记：郑经死期和继承结果可靠；宫廷政变过程主要依郑氏后出材料与清方情报

\paragraph{康熙二十年（1681年）：郑经去世，东宁发生继承冲突，郑克塽最终掌权}
\eventanchor{EQ-1681-ZHENG-JING-DEATH-DECISION}{郑经去世，东宁发生继承冲突，郑克塽最终掌权} 标记本节点定位围绕「郑经去世，东宁发生继承冲突，郑克塽最终掌权」形成的处置与调度记录，涉及东宁郑氏政权。京师与多地军政线路相接的尺度不是抽象背景：郑经长期在大陆作战并承受台湾财政压力，其死亡使宗族、将领和文官围绕继承重组。《清圣祖实录》康熙二十年台湾情报；《清史稿·郑成功传附郑经》；《台湾外纪》与荷兰商馆残存通信所见的顺序随后指向东宁短期政治不稳，清廷得以评估以水师进攻台湾的风险和机会，但原有置信注记：郑经死期和继承结果可靠；宫廷政变过程主要依郑氏后出材料与清方情报

\paragraph{康熙二十年（1681年）：郑经去世，东宁发生继承冲突，郑克塽最终掌权}
在京师与多地军政线路相接的尺度，\eventanchor{EQ-1681-ZHENG-JING-DEATH-PRELUDE}{郑经去世，东宁发生继承冲突，郑克塽最终掌权} 让东宁郑氏政权的一个选择或压力有了可查的坐标。本节点回看「郑经去世，东宁发生继承冲突，郑克塽最终掌权」之前已可见的条件。前因可写为郑经长期在大陆作战并承受台湾财政压力，其死亡使宗族、将领和文官围绕继承重组；《清圣祖实录》康熙二十年台湾情报；《清史稿·郑成功传附郑经》；《台湾外纪》与荷兰商馆残存通信限定了记录，后续变化是东宁短期政治不稳，清廷得以评估以水师进攻台湾的风险和机会。原有置信注记：郑经死期和继承结果可靠；宫廷政变过程主要依郑氏后出材料与清方情报

\paragraph{康熙二十一年（1682年）：清廷任命施琅统率福建水师，筹备进攻台湾}
\eventanchor{EQ-1682-SHI-LANG}{清廷任命施琅统率福建水师，筹备进攻台湾} 所关联的是清与东宁郑氏政权在港口、海峡与跨海航路相连的尺度中的行动。本节点叙述该事件本身。它从三藩平定和东宁继承危机减少内陆牵制，厦金旧将对航道与郑氏军制的了解提供条件展开，借《清圣祖实录》康熙二十一年福建海防条；《清史稿·施琅传》；《康熙朝满文朱批奏折全译》康熙二十一年台湾军务奏折进入可检验的年序，并以清廷从迁界防御转向跨海进攻，水师训练、粮械和季风仍决定行动可行性影响下一段安排。原有置信注记：任命和筹备可靠；施琅战略构想与朝廷内部争论须以奏折、朱批和后出传记区别处理

\paragraph{康熙二十一年（1682年）：清廷任命施琅统率福建水师，筹备进攻台湾}
港口、海峡与跨海航路相连的尺度中的\eventanchor{EQ-1682-SHI-LANG-AFTERMATH}{清廷任命施琅统率福建水师，筹备进攻台湾} 不能离开清与东宁郑氏政权来解释。本节点追踪「清廷任命施琅统率福建水师，筹备进攻台湾」之后可见的余波。三藩平定和东宁继承危机减少内陆牵制，厦金旧将对航道与郑氏军制的了解提供条件说明其形成的条件；《清圣祖实录》康熙二十一年福建海防条；《清史稿·施琅传》；《康熙朝满文朱批奏折全译》康熙二十一年台湾军务奏折提供来源路径，清廷从迁界防御转向跨海进攻，水师训练、粮械和季风仍决定行动可行性则是材料能够支持的近时结果。原有置信注记：任命和筹备可靠；施琅战略构想与朝廷内部争论须以奏折、朱批和后出传记区别处理

\paragraph{康熙二十一年（1682年）：清廷任命施琅统率福建水师，筹备进攻台湾}
\eventanchor{EQ-1682-SHI-LANG-DECISION}{清廷任命施琅统率福建水师，筹备进攻台湾} 把清与东宁郑氏政权放回港口、海峡与跨海航路相连的尺度的道路、城镇、边界或制度网络。本节点定位围绕「清廷任命施琅统率福建水师，筹备进攻台湾」形成的处置与调度记录。三藩平定和东宁继承危机减少内陆牵制，厦金旧将对航道与郑氏军制的了解提供条件。本段采用《清圣祖实录》康熙二十一年福建海防条；《清史稿·施琅传》；《康熙朝满文朱批奏折全译》康熙二十一年台湾军务奏折核对其顺序，随后可追到清廷从迁界防御转向跨海进攻，水师训练、粮械和季风仍决定行动可行性。原有置信注记：任命和筹备可靠；施琅战略构想与朝廷内部争论须以奏折、朱批和后出传记区别处理

\paragraph{康熙二十一年（1682年）：清廷任命施琅统率福建水师，筹备进攻台湾}
本节点回看「清廷任命施琅统率福建水师，筹备进攻台湾」之前已可见的条件，故\eventanchor{EQ-1682-SHI-LANG-PRELUDE}{清廷任命施琅统率福建水师，筹备进攻台湾} 不只是一个年份标签。清与东宁郑氏政权在港口、海峡与跨海航路相连的尺度承受的压力来自三藩平定和东宁继承危机减少内陆牵制，厦金旧将对航道与郑氏军制的了解提供条件。《清圣祖实录》康熙二十一年福建海防条；《清史稿·施琅传》；《康熙朝满文朱批奏折全译》康熙二十一年台湾军务奏折可回查该节点；清廷从迁界防御转向跨海进攻，水师训练、粮械和季风仍决定行动可行性是其进入后续年序的方式。原有置信注记：任命和筹备可靠；施琅战略构想与朝廷内部争论须以奏折、朱批和后出传记区别处理

\paragraph{康熙二十二年六月（1683年）：施琅水师在澎湖击败郑氏舰队，取得通往台湾本岛的制海优势}
\eventanchor{EQ-1683-PENGHU}{施琅水师在澎湖击败郑氏舰队，取得通往台湾本岛的制海优势} 所见的行动者为清与东宁郑氏政权，空间尺度为港口、海峡与跨海航路相连的尺度。本节点叙述该事件本身。清军利用季风窗口、福建补给和施琅对郑氏舰队的熟悉发动集中攻击，郑氏难以同时守澎湖与本岛使它成为具体事件链的一环；《清圣祖实录》康熙二十二年六月；《清史稿·施琅传》；施琅《飞报澎湖大捷疏》与台湾地方志与郑氏失去海上屏障，台湾方面转入投降谈判；澎湖伤亡和战术细节仍需多源检验。原有置信注记：澎湖战役结果可靠；舰船、兵员和死伤数字多来自胜方战报，须以郑氏和海域材料约束

\paragraph{康熙二十二年六月（1683年）：施琅水师在澎湖击败郑氏舰队，取得通往台湾本岛的制海优势}
从港口、海峡与跨海航路相连的尺度看，\eventanchor{EQ-1683-PENGHU-AFTERMATH}{施琅水师在澎湖击败郑氏舰队，取得通往台湾本岛的制海优势} 留下本节点追踪「施琅水师在澎湖击败郑氏舰队，取得通往台湾本岛的制海优势」之后可见的余波的材料坐标。清与东宁郑氏政权所面对的条件是清军利用季风窗口、福建补给和施琅对郑氏舰队的熟悉发动集中攻击，郑氏难以同时守澎湖与本岛。《清圣祖实录》康熙二十二年六月；《清史稿·施琅传》；施琅《飞报澎湖大捷疏》与台湾地方志支持的结果为郑氏失去海上屏障，台湾方面转入投降谈判；澎湖伤亡和战术细节仍需多源检验。原有置信注记：澎湖战役结果可靠；舰船、兵员和死伤数字多来自胜方战报，须以郑氏和海域材料约束

\paragraph{康熙二十二年六月（1683年）：施琅水师在澎湖击败郑氏舰队，取得通往台湾本岛的制海优势}
\eventanchor{EQ-1683-PENGHU-DECISION}{施琅水师在澎湖击败郑氏舰队，取得通往台湾本岛的制海优势} 记录清与东宁郑氏政权在港口、海峡与跨海航路相连的尺度的一次变化。本节点定位围绕「施琅水师在澎湖击败郑氏舰队，取得通往台湾本岛的制海优势」形成的处置与调度记录。其发生与清军利用季风窗口、福建补给和施琅对郑氏舰队的熟悉发动集中攻击，郑氏难以同时守澎湖与本岛相连；《清圣祖实录》康熙二十二年六月；《清史稿·施琅传》；施琅《飞报澎湖大捷疏》与台湾地方志把事件系回来源，郑氏失去海上屏障，台湾方面转入投降谈判；澎湖伤亡和战术细节仍需多源检验显示压力如何向后转移。原有置信注记：澎湖战役结果可靠；舰船、兵员和死伤数字多来自胜方战报，须以郑氏和海域材料约束

\paragraph{康熙二十二年六月（1683年）：施琅水师在澎湖击败郑氏舰队，取得通往台湾本岛的制海优势}
\eventanchor{EQ-1683-PENGHU-PRELUDE}{施琅水师在澎湖击败郑氏舰队，取得通往台湾本岛的制海优势} 的地点与行动范围属于港口、海峡与跨海航路相连的尺度，行动者是清与东宁郑氏政权。本节点回看「施琅水师在澎湖击败郑氏舰队，取得通往台湾本岛的制海优势」之前已可见的条件。清军利用季风窗口、福建补给和施琅对郑氏舰队的熟悉发动集中攻击，郑氏难以同时守澎湖与本岛。读《清圣祖实录》康熙二十二年六月；《清史稿·施琅传》；施琅《飞报澎湖大捷疏》与台湾地方志时可见其记录边界，最小后果只能写作郑氏失去海上屏障，台湾方面转入投降谈判；澎湖伤亡和战术细节仍需多源检验。原有置信注记：澎湖战役结果可靠；舰船、兵员和死伤数字多来自胜方战报，须以郑氏和海域材料约束

\paragraph{康熙二十二年八月（1683年）：郑克塽向清廷投降，东宁政权结束}
\eventanchor{EQ-1683-TAIWAN-SURRENDER}{郑克塽向清廷投降，东宁政权结束} 是清与东宁郑氏政权、台湾社会在港口、海峡与跨海航路相连的尺度留下的编年标记。本节点叙述该事件本身。它从澎湖失败切断东宁海防与外援预期，继承危机使郑氏难以继续集中抵抗走来；《清圣祖实录》康熙二十二年八月；《清史稿·郑成功传附郑克塽》；施琅奏折与《台湾府志》固定可证部分，清廷取得台湾政治接收权，后续设治、迁界调整和土地社会重组并非投降即告完成构成下一步的可见结果。原有置信注记：投降与清廷接收主干可靠；投降条件、官兵安置和原住民、移民社会变化须分层记录

\paragraph{康熙二十二年八月（1683年）：郑克塽向清廷投降，东宁政权结束}
本节点追踪「郑克塽向清廷投降，东宁政权结束」之后可见的余波：\eventanchor{EQ-1683-TAIWAN-SURRENDER-AFTERMATH}{郑克塽向清廷投降，东宁政权结束}。清与东宁郑氏政权、台湾社会在港口、海峡与跨海航路相连的尺度的处境可由澎湖失败切断东宁海防与外援预期，继承危机使郑氏难以继续集中抵抗说明。《清圣祖实录》康熙二十二年八月；《清史稿·郑成功传附郑克塽》；施琅奏折与《台湾府志》是本段材料入口，因而只能据此追踪到清廷取得台湾政治接收权，后续设治、迁界调整和土地社会重组并非投降即告完成。原有置信注记：投降与清廷接收主干可靠；投降条件、官兵安置和原住民、移民社会变化须分层记录

\paragraph{康熙二十二年八月（1683年）：郑克塽向清廷投降，东宁政权结束}
\eventanchor{EQ-1683-TAIWAN-SURRENDER-DECISION}{郑克塽向清廷投降，东宁政权结束} 发生在港口、海峡与跨海航路相连的尺度，牵动清与东宁郑氏政权、台湾社会。本节点定位围绕「郑克塽向清廷投降，东宁政权结束」形成的处置与调度记录。澎湖失败切断东宁海防与外援预期，继承危机使郑氏难以继续集中抵抗是其结构性条件；《清圣祖实录》康熙二十二年八月；《清史稿·郑成功传附郑克塽》；施琅奏折与《台湾府志》留下的证据把读者带到清廷取得台湾政治接收权，后续设治、迁界调整和土地社会重组并非投降即告完成。原有置信注记：投降与清廷接收主干可靠；投降条件、官兵安置和原住民、移民社会变化须分层记录

\paragraph{康熙二十二年八月（1683年）：郑克塽向清廷投降，东宁政权结束}
\eventanchor{EQ-1683-TAIWAN-SURRENDER-PRELUDE}{郑克塽向清廷投降，东宁政权结束} 对应本节点回看「郑克塽向清廷投降，东宁政权结束」之前已可见的条件。在港口、海峡与跨海航路相连的尺度，清与东宁郑氏政权、台湾社会无法绕开澎湖失败切断东宁海防与外援预期，继承危机使郑氏难以继续集中抵抗。依据《清圣祖实录》康熙二十二年八月；《清史稿·郑成功传附郑克塽》；施琅奏折与《台湾府志》，本段把后续影响限定为清廷取得台湾政治接收权，后续设治、迁界调整和土地社会重组并非投降即告完成。原有置信注记：投降与清廷接收主干可靠；投降条件、官兵安置和原住民、移民社会变化须分层记录

\paragraph{康熙二十三年（1684年）：清廷调整迁界与海禁政策，允许沿海居民复界并逐步恢复民间海贸}
\eventanchor{EQ-1684-SEA-BAN-LIFT}{清廷调整迁界与海禁政策，允许沿海居民复界并逐步恢复民间海贸}。本节点叙述该事件本身。清与闽粤沿海在港口、海峡与跨海航路相连的尺度的处境首先受台湾投降降低清廷对郑氏补给网络的担忧，长期迁界的税源、民生和治安代价亦需处理制约。《清圣祖实录》康熙二十三年复界海禁条；《清史稿·食货志》；《福建通志》《广东通志》沿海复界条留下可核的线索；可辨的后果是沿海耕作、渔盐和航运开始恢复，但贸易许可、海防和地方执行塑造不同港口开放节奏。原有置信注记：政策转向可靠；复界日期、港口开放和走私变化因省县而异，不能由中央谕旨推全国

\paragraph{康熙二十三年（1684年）：清廷调整迁界与海禁政策，允许沿海居民复界并逐步恢复民间海贸}
\eventanchor{EQ-1684-SEA-BAN-LIFT-AFTERMATH}{清廷调整迁界与海禁政策，允许沿海居民复界并逐步恢复民间海贸} 置于港口、海峡与跨海航路相连的尺度中阅读，本节点追踪「清廷调整迁界与海禁政策，允许沿海居民复界并逐步恢复民间海贸」之后可见的余波。台湾投降降低清廷对郑氏补给网络的担忧，长期迁界的税源、民生和治安代价亦需处理构成行动者清与闽粤沿海必须面对的条件。取证从《清圣祖实录》康熙二十三年复界海禁条；《清史稿·食货志》；《福建通志》《广东通志》沿海复界条出发，因而只能把变化谨慎地连到沿海耕作、渔盐和航运开始恢复，但贸易许可、海防和地方执行塑造不同港口开放节奏。原有置信注记：政策转向可靠；复界日期、港口开放和走私变化因省县而异，不能由中央谕旨推全国

\paragraph{康熙二十三年（1684年）：清廷调整迁界与海禁政策，允许沿海居民复界并逐步恢复民间海贸}
\eventanchor{EQ-1684-SEA-BAN-LIFT-DECISION}{清廷调整迁界与海禁政策，允许沿海居民复界并逐步恢复民间海贸} 标记本节点定位围绕「清廷调整迁界与海禁政策，允许沿海居民复界并逐步恢复民间海贸」形成的处置与调度记录，涉及清与闽粤沿海。港口、海峡与跨海航路相连的尺度不是抽象背景：台湾投降降低清廷对郑氏补给网络的担忧，长期迁界的税源、民生和治安代价亦需处理。《清圣祖实录》康熙二十三年复界海禁条；《清史稿·食货志》；《福建通志》《广东通志》沿海复界条所见的顺序随后指向沿海耕作、渔盐和航运开始恢复，但贸易许可、海防和地方执行塑造不同港口开放节奏，但原有置信注记：政策转向可靠；复界日期、港口开放和走私变化因省县而异，不能由中央谕旨推全国

\paragraph{康熙二十三年（1684年）：清廷调整迁界与海禁政策，允许沿海居民复界并逐步恢复民间海贸}
在港口、海峡与跨海航路相连的尺度，\eventanchor{EQ-1684-SEA-BAN-LIFT-PRELUDE}{清廷调整迁界与海禁政策，允许沿海居民复界并逐步恢复民间海贸} 让清与闽粤沿海的一个选择或压力有了可查的坐标。本节点回看「清廷调整迁界与海禁政策，允许沿海居民复界并逐步恢复民间海贸」之前已可见的条件。前因可写为台湾投降降低清廷对郑氏补给网络的担忧，长期迁界的税源、民生和治安代价亦需处理；《清圣祖实录》康熙二十三年复界海禁条；《清史稿·食货志》；《福建通志》《广东通志》沿海复界条限定了记录，后续变化是沿海耕作、渔盐和航运开始恢复，但贸易许可、海防和地方执行塑造不同港口开放节奏。原有置信注记：政策转向可靠；复界日期、港口开放和走私变化因省县而异，不能由中央谕旨推全国

\paragraph{康熙二十三年（1684年）：清廷设置台湾府及诸县，隶属福建省，开始以府县和驻防制度经营台湾}
\eventanchor{EQ-1684-TAIWAN-PREFECTURE}{清廷设置台湾府及诸县，隶属福建省，开始以府县和驻防制度经营台湾} 所关联的是清与台湾在港口、海峡与跨海航路相连的尺度中的行动。本节点叙述该事件本身。它从东宁投降后清廷须在弃守与设治间权衡海防、税源、移民和驻军成本展开，借《清圣祖实录》康熙二十三年台湾设治条；《清史稿·地理志》；《台湾府志》与福建省档案相关奏议进入可检验的年序，并以台湾纳入福建行政体系，但沿海、平原、内山和原住民社群与官府关系呈现不均衡接触影响下一段安排。原有置信注记：设府设县可靠；机构设置不等于全岛有效治理，原住民地区和内山控制尤须保留

\paragraph{康熙二十三年（1684年）：清廷设置台湾府及诸县，隶属福建省，开始以府县和驻防制度经营台湾}
港口、海峡与跨海航路相连的尺度中的\eventanchor{EQ-1684-TAIWAN-PREFECTURE-AFTERMATH}{清廷设置台湾府及诸县，隶属福建省，开始以府县和驻防制度经营台湾} 不能离开清与台湾来解释。本节点追踪「清廷设置台湾府及诸县，隶属福建省，开始以府县和驻防制度经营台湾」之后可见的余波。东宁投降后清廷须在弃守与设治间权衡海防、税源、移民和驻军成本说明其形成的条件；《清圣祖实录》康熙二十三年台湾设治条；《清史稿·地理志》；《台湾府志》与福建省档案相关奏议提供来源路径，台湾纳入福建行政体系，但沿海、平原、内山和原住民社群与官府关系呈现不均衡接触则是材料能够支持的近时结果。原有置信注记：设府设县可靠；机构设置不等于全岛有效治理，原住民地区和内山控制尤须保留

\paragraph{康熙二十三年（1684年）：清廷设置台湾府及诸县，隶属福建省，开始以府县和驻防制度经营台湾}
\eventanchor{EQ-1684-TAIWAN-PREFECTURE-DECISION}{清廷设置台湾府及诸县，隶属福建省，开始以府县和驻防制度经营台湾} 把清与台湾放回港口、海峡与跨海航路相连的尺度的道路、城镇、边界或制度网络。本节点定位围绕「清廷设置台湾府及诸县，隶属福建省，开始以府县和驻防制度经营台湾」形成的处置与调度记录。东宁投降后清廷须在弃守与设治间权衡海防、税源、移民和驻军成本。本段采用《清圣祖实录》康熙二十三年台湾设治条；《清史稿·地理志》；《台湾府志》与福建省档案相关奏议核对其顺序，随后可追到台湾纳入福建行政体系，但沿海、平原、内山和原住民社群与官府关系呈现不均衡接触。原有置信注记：设府设县可靠；机构设置不等于全岛有效治理，原住民地区和内山控制尤须保留

\paragraph{康熙二十三年（1684年）：清廷设置台湾府及诸县，隶属福建省，开始以府县和驻防制度经营台湾}
本节点回看「清廷设置台湾府及诸县，隶属福建省，开始以府县和驻防制度经营台湾」之前已可见的条件，故\eventanchor{EQ-1684-TAIWAN-PREFECTURE-PRELUDE}{清廷设置台湾府及诸县，隶属福建省，开始以府县和驻防制度经营台湾} 不只是一个年份标签。清与台湾在港口、海峡与跨海航路相连的尺度承受的压力来自东宁投降后清廷须在弃守与设治间权衡海防、税源、移民和驻军成本。《清圣祖实录》康熙二十三年台湾设治条；《清史稿·地理志》；《台湾府志》与福建省档案相关奏议可回查该节点；台湾纳入福建行政体系，但沿海、平原、内山和原住民社群与官府关系呈现不均衡接触是其进入后续年序的方式。原有置信注记：设府设县可靠；机构设置不等于全岛有效治理，原住民地区和内山控制尤须保留

\paragraph{康熙六十年（1721年）：朱一贵在台湾发动反清起事，一度攻入府城，清廷随后派兵恢复控制}
\eventanchor{EQ-1721-ZHU-YIGUI}{朱一贵在台湾发动反清起事，一度攻入府城，清廷随后派兵恢复控制} 所见的行动者为清与台湾朱一贵集团，空间尺度为港口、海峡与跨海航路相连的尺度。本节点叙述该事件本身。台湾设治后垦殖、租佃、族群关系与地方行政摩擦累积，官府征收和治安处置成为触发因素使它成为具体事件链的一环；《清圣祖实录》康熙六十年台湾军报；《清史稿·朱一贵传》；《台湾府志》、福建督抚奏折与契约文书与清廷暴露跨海调兵与府县治理薄弱环节，后续强化驻防和行政整顿；不能将起事简化为单一族群反应。原有置信注记：起事、府城易手与镇压主干可靠；参与群体、赋税原因和伤亡数字受官府“逆匪”分类影响

\paragraph{康熙六十年（1721年）：朱一贵在台湾发动反清起事，一度攻入府城，清廷随后派兵恢复控制}
从港口、海峡与跨海航路相连的尺度看，\eventanchor{EQ-1721-ZHU-YIGUI-AFTERMATH}{朱一贵在台湾发动反清起事，一度攻入府城，清廷随后派兵恢复控制} 留下本节点追踪「朱一贵在台湾发动反清起事，一度攻入府城，清廷随后派兵恢复控制」之后可见的余波的材料坐标。清与台湾朱一贵集团所面对的条件是台湾设治后垦殖、租佃、族群关系与地方行政摩擦累积，官府征收和治安处置成为触发因素。《清圣祖实录》康熙六十年台湾军报；《清史稿·朱一贵传》；《台湾府志》、福建督抚奏折与契约文书支持的结果为清廷暴露跨海调兵与府县治理薄弱环节，后续强化驻防和行政整顿；不能将起事简化为单一族群反应。原有置信注记：起事、府城易手与镇压主干可靠；参与群体、赋税原因和伤亡数字受官府“逆匪”分类影响

\paragraph{康熙六十年（1721年）：朱一贵在台湾发动反清起事，一度攻入府城，清廷随后派兵恢复控制}
\eventanchor{EQ-1721-ZHU-YIGUI-DECISION}{朱一贵在台湾发动反清起事，一度攻入府城，清廷随后派兵恢复控制} 记录清与台湾朱一贵集团在港口、海峡与跨海航路相连的尺度的一次变化。本节点定位围绕「朱一贵在台湾发动反清起事，一度攻入府城，清廷随后派兵恢复控制」形成的处置与调度记录。其发生与台湾设治后垦殖、租佃、族群关系与地方行政摩擦累积，官府征收和治安处置成为触发因素相连；《清圣祖实录》康熙六十年台湾军报；《清史稿·朱一贵传》；《台湾府志》、福建督抚奏折与契约文书把事件系回来源，清廷暴露跨海调兵与府县治理薄弱环节，后续强化驻防和行政整顿；不能将起事简化为单一族群反应显示压力如何向后转移。原有置信注记：起事、府城易手与镇压主干可靠；参与群体、赋税原因和伤亡数字受官府“逆匪”分类影响

\paragraph{康熙六十年（1721年）：朱一贵在台湾发动反清起事，一度攻入府城，清廷随后派兵恢复控制}
\eventanchor{EQ-1721-ZHU-YIGUI-PRELUDE}{朱一贵在台湾发动反清起事，一度攻入府城，清廷随后派兵恢复控制} 的地点与行动范围属于港口、海峡与跨海航路相连的尺度，行动者是清与台湾朱一贵集团。本节点回看「朱一贵在台湾发动反清起事，一度攻入府城，清廷随后派兵恢复控制」之前已可见的条件。台湾设治后垦殖、租佃、族群关系与地方行政摩擦累积，官府征收和治安处置成为触发因素。读《清圣祖实录》康熙六十年台湾军报；《清史稿·朱一贵传》；《台湾府志》、福建督抚奏折与契约文书时可见其记录边界，最小后果只能写作清廷暴露跨海调兵与府县治理薄弱环节，后续强化驻防和行政整顿；不能将起事简化为单一族群反应。原有置信注记：起事、府城易手与镇压主干可靠；参与群体、赋税原因和伤亡数字受官府“逆匪”分类影响